\scanpage{39}
\deleted{\unsure{Soit}}

Soit $K$ un corps topologique isomorphe à $\mathbb C$, $X$ un
schéma algébrique défini sur $K$. Alors $X_{\mathrm{top}}$ est défini,
et aussi $H^*(X_{\mathrm{top}},K)$. D'autre part, \inserted{si $X$ est non sing.,} la cohomologie
de Hodge est définie, soit $H^*(X)$ [indépendant de
la structure topologique de $K$]. Alors le choix
d'un isom top. $u:K\unsure{-|\to}\mathbb C$ définit un isomorphisme
\[
H^*(X_{\mathrm{top}},K)\xrightarrow{\sim} H^*(X)
\]
et on vérifie aussitôt que cet isomorphisme ne
dépend pas du choix de $u$, [En d'autres termes,
si $K=\mathbb C$, alors les deux homomorphismes
\diagram{\begin{tikzcd}[ampersand replacement=\&,row sep=small]
\& H^*(X) \arrow[dd,"\simeq\ (\text{antilinéaire})"]\\
H^*(X_{\mathrm{top}},\mathbb C) \arrow[ur] \arrow[dr] \arrow[r,phantom,"\text{non comm.}" description] \& \\
\& H^*(X^c)
\end{tikzcd}}
(où $c:\mathbb C-\mathbb C$ est la conjugaison complexe)
sont \deleted{conjugués} déduits
l'un de l'autre par application de $c$ aux coefficients de $H^*(X_{\mathrm{top}},\mathbb C)$]

\inserted{Cela nous permet, notant qu'on a un \unsure{isom|hom|homom.} injectif continu unique
$\unsure{u|i}:\mathbb R\to\mathbb C$, de définir $H^*(X_{\mathrm{top}},\mathbb R)\hookrightarrow H^*(X)$ (même si $X$ pas définissable
% margin note written sideways
sur $\mathbb R$ !) d'où un sous-\inserted{$\mathbb R$-}espace
de $H^*(X)$, \deleted{défini sur $\mathbb R$, \unsure{et}} noté $H^*(X_{\unsure{\tau|+|?}},\mathbb R)$
\deleted{$\unsure{H^*(X)^+}$}, ou \deleted{\unsure{cet} \unsure{dépend de la topologie de} $K$},
\unsure{aut.nt|anti} $H^*(X)^-=\sqrt{-1}H^*(X)^+$.}

Supposons $X$ projective non singulière, alors
le choix d'un isom top. $u:K\xrightarrow{\sim}\mathbb C$ définit
une décomposition
\[
\varphi:\coprod_{p,q}H^q(X,\Omega_X^p)\xrightarrow{\sim}H^*(X)
\]
compatible avec la filtration de $H^*(X)$. On
vérifie encore que $\varphi$ est indépendant du choix de $u$.
% archive sheet number: 38
